\documentclass[a4paper,12pt]{extarticle}

\usepackage{fontspec}
\usepackage{polyglossia}
\setdefaultlanguage{russian}
\setotherlanguage{english}
\setmainfont{Times New Roman}
\setmonofont{Courier New}
\newfontfamily\cyrillicfonttt{Courier New}

\usepackage[nodisplayskipstretch]{setspace}
\onehalfspacing

\usepackage[a4paper, left=30mm, right=15mm, top=20mm, bottom=20mm]{geometry}
\setlength{\parindent}{1.25cm}
\setlength{\parskip}{0pt}
\usepackage{indentfirst}

\usepackage{fancyhdr}
\pagestyle{fancy}
\fancyhf{}
\fancyfoot[C]{\thepage}
\renewcommand{\headrulewidth}{0pt}
\renewcommand{\footrulewidth}{0pt}
\fancypagestyle{plain}{
  \fancyhf{}
  \fancyfoot[C]{\thepage}
  \renewcommand{\headrulewidth}{0pt}
  \renewcommand{\footrulewidth}{0pt}
}

\usepackage{titlesec}
\titleformat{\section}[block]
  {\bfseries\fontsize{14pt}{14pt}\selectfont}
  {\thesection}
  {1em}
  {}
\titlespacing*{\section}{\parindent}{3\baselineskip}{3\baselineskip}

\titleformat{\subsection}[block]
  {\bfseries\normalsize}
  {\thesubsection}
  {1em}
  {}
\titlespacing*{\subsection}{\parindent}{\baselineskip}{\baselineskip}

\titleformat{\subsubsection}[block]
  {\bfseries\normalsize}
  {\thesubsubsection}
  {1em}
  {}
\titlespacing*{\subsubsection}{\parindent}{\baselineskip}{\baselineskip}

\usepackage{tocloft}
\renewcommand{\contentsname}{СОДЕРЖАНИЕ}
\renewcommand{\cfttoctitlefont}{\hfill\bfseries}
\renewcommand{\cftaftertoctitle}{\hfill}
\setlength{\cftbeforetoctitleskip}{0pt}
\setlength{\cftaftertoctitleskip}{2\baselineskip}
\renewcommand{\cftsecleader}{\cftdotfill{\cftdotsep}}

\usepackage{amsmath}
\usepackage{amssymb}
\usepackage{graphicx}
\usepackage{pdfpages}
\usepackage{array}
\usepackage{booktabs}
\usepackage{longtable}
\usepackage{float}

\usepackage{caption}
\DeclareCaptionLabelSeparator{emdash}{~--~}
\captionsetup[figure]{labelsep=emdash, justification=centering, singlelinecheck=false}
\captionsetup[table]{labelsep=emdash, justification=raggedright, singlelinecheck=false, position=top}

\usepackage{enumitem}
\setlist[itemize]{label=--, itemsep=0pt, topsep=0pt}
\setlist[enumerate]{itemsep=0pt, topsep=0pt}

\usepackage{listings}
\usepackage{xcolor}
\lstset{
  basicstyle=\rmfamily\small,
  breaklines=true,
  frame=single,
  columns=fullflexible,
  keepspaces=true,
  showstringspaces=false,
  xleftmargin=\parindent
}

\usepackage{hyperref}
\hypersetup{
  unicode=true,
  colorlinks=false,
  pdfborder={0 0 0}
}
\urlstyle{same}

\makeatletter
\renewcommand\@biblabel[1]{#1.}
\renewenvironment{thebibliography}[1]
     {\list{\@biblabel{\@arabic\c@enumiv}}%
           {\settowidth\labelwidth{\@biblabel{#1}}%
            \leftmargin\labelwidth
            \advance\leftmargin\labelsep
            \setlength{\itemsep}{0pt}%
            \setlength{\parsep}{0pt}%
            \setlength{\parskip}{0pt}%
            \setlength{\topsep}{0pt}%
            \usecounter{enumiv}%
            \let\p@enumiv\@empty
            \renewcommand\theenumiv{\@arabic\c@enumiv}}%
      \sloppy
      \clubpenalty4000
      \widowpenalty4000
      \sfcode`\.\@m}
     {\def\@noitemerr{\@latex@warning{Empty `thebibliography' environment}}%
      \endlist}
\makeatother

\newcommand{\structuralsection}[1]{%
  \clearpage
  \phantomsection
  \addcontentsline{toc}{section}{#1}%
  \begin{center}
    \textbf{\MakeUppercase{#1}}
  \end{center}
}

\newcommand{\appendixsection}[2]{%
  \clearpage
  \phantomsection
  \addcontentsline{toc}{section}{Приложение #1. #2}%
  \begin{flushright}
    \textbf{Приложение #1}
  \end{flushright}
  \begin{center}
    \textbf{#2}
  \end{center}
}

\sloppy


\newcommand{\EducationProgram}{[код и наименование образовательной программы]}
\newcommand{\TrainingDirection}{[код и наименование направления подготовки]}
\newcommand{\StudentName}{[Фамилия Имя Отчество]}
\newcommand{\GroupNumber}{[номер группы]}
\newcommand{\CompanySupervisor}{[ФИО, должность, организация]}
\newcommand{\UniversitySupervisor}{[ФИО, должность]}
\newcommand{\PracticePeriod}{[сроки прохождения практики]}
\newcommand{\PracticePlace}{Университет ИТМО, система дистанционного обучения}
\newcommand{\PracticeTopic}{Разработка интерактивного тренажёра <<Алгоритм Флёри: построение эйлерова цикла в неориентированном графе>> для системы дистанционного обучения ИТМО}
\newcommand{\PracticeYear}{2026}

\begin{document}

\iffalse
\begin{titlepage}
\thispagestyle{empty}
\begin{center}
Министерство науки и высшего образования Российской Федерации\\
\textbf{ФЕДЕРАЛЬНОЕ ГОСУДАРСТВЕННОЕ АВТОНОМНОЕ ОБРАЗОВАТЕЛЬНОЕ}\\
\textbf{УЧРЕЖДЕНИЕ ВЫСШЕГО ОБРАЗОВАНИЯ}\\
\textbf{<<НАЦИОНАЛЬНЫЙ ИССЛЕДОВАТЕЛЬСКИЙ УНИВЕРСИТЕТ ИТМО>>}\\
\textbf{(Университет ИТМО)}\\[8pt]
\textbf{Факультет программной инженерии и компьютерной техники}
\end{center}

\vspace{8pt}
\noindent Образовательная программа:\\[4pt]
\underline{\makebox[\textwidth][l]{\EducationProgram}}\\[8pt]
\noindent Направление подготовки (специальность):\\[4pt]
\underline{\makebox[\textwidth][l]{\TrainingDirection}}\\[20pt]

\begin{center}
{\large \textbf{ОТЧЁТ}}\\[6pt]
о производственной, технологической\\
(проектно-технологической) практике
\end{center}

\vspace{12pt}
\noindent Тема задания:\\[4pt]
\begin{center}
\parbox{0.92\textwidth}{\centering \PracticeTopic}
\end{center}
\vspace{8pt}
\noindent Обучающийся \underline{\makebox[11.7cm][l]{\StudentName}}\\[8pt]
\noindent Номер группы \underline{\makebox[4cm][l]{\GroupNumber}}\\[12pt]
\noindent Руководитель практики от профильной организации: \\
\underline{\makebox[15.2cm][l]{\CompanySupervisor}}\\[12pt]
\noindent Руководитель практики от Университета: \\
\underline{\makebox[15.2cm][l]{\UniversitySupervisor}}\\[12pt]
\noindent Место прохождения практики: \PracticePlace\\[8pt]
\noindent Сроки прохождения практики: \PracticePeriod\\[14pt]
\noindent Практика пройдена с оценкой \underline{\makebox[3cm][l]{}}\\[12pt]
\noindent Дата \underline{\makebox[3cm][l]{}}\\

\vfill
\begin{center}
Санкт-Петербург\\
\PracticeYear{} г.
\end{center}
\end{titlepage}
\fi
\includepdf[pages=1,pagecommand={\thispagestyle{empty}}]{Титульный лист.pdf}
\setcounter{page}{2}

\tableofcontents

\structuralsection{Обозначения и сокращения}

\noindent
API -- application programming interface, программный интерфейс приложения.\\
CSS -- cascading style sheets, каскадные таблицы стилей.\\
HTML -- hypertext markup language, язык гипертекстовой разметки.\\
ITMO -- Университет ИТМО.\\
JSON -- javascript object notation, формат обмена структурированными данными.\\
LMS -- learning management system, система дистанционного обучения.\\
UI -- user interface, пользовательский интерфейс.

\structuralsection{Характеристика организации (места прохождения практики)}

Местом прохождения практики являлся Университет ИТМО, а предметной средой работ выступала система дистанционного обучения, в которой размещаются интерактивные задания для студентов. Такая среда предъявляет к программным средствам одновременно методические и технические требования: интерфейс должен быть понятным для обучающихся, алгоритмическая логика -- корректной, а развёртывание -- воспроизводимым в изолированной контейнерной инфраструктуре.

Практическая работа выполнялась в контексте разработки учебного веб-приложения, предназначенного для закрепления темы эйлеровых циклов и алгоритма Флёри. Целевой аудиторией тренажёра являются студенты, изучающие дискретную математику и теорию графов, а также преподаватели, которым необходим инструмент для выдачи самостоятельных заданий, автоматизированной проверки и демонстрации корректного построения эйлерова цикла.

Особенностью места прохождения практики является ориентация на дистанционный формат. Это означает, что программный продукт должен одинаково поддерживать учебный сценарий, режим контролируемого тестирования, передачу результатов в внешнюю среду и запуск в стандартной инфраструктуре LMS. Подобные требования соответствуют целям производственной, технологической (проектно-технологической) практики, в рамках которой результатом работы является не только код, но и законченный цифровой артефакт, пригодный для эксплуатации в учебном процессе~\cite{markina2020,gost732}.

\structuralsection{Введение}

Производственная, технологическая (проектно-технологическая) практика направлена на получение и закрепление навыков практической профессиональной деятельности, связанных с анализом требований, проектированием, программной реализацией и документированием программных средств~\cite{markina2020}. Для ИТМО особенно актуальны разработки, которые могут быть интегрированы в цифровую образовательную среду и использоваться в дистанционном обучении.

Тема настоящей практики -- разработка интерактивного тренажёра <<Алгоритм Флёри: построение эйлерова цикла в неориентированном графе>> для системы дистанционного обучения ИТМО. Выбор темы обусловлен двумя факторами. Во-первых, алгоритм Флёри является классическим алгоритмом теории графов и используется для объяснения условий существования и построения эйлеровых циклов~\cite{wilson2010,west2001}. Во-вторых, для эффективного освоения этой темы студентам недостаточно статического текста: необходимы визуализация графа, пошаговое прохождение и проверка собственных действий.

Объектом практики является образовательное веб-приложение, предназначенное для работы с неориентированными графами. Предмет практики -- методы проектирования и реализации интерактивного тренажёра, обеспечивающего генерацию эйлеровых графов, демонстрацию алгоритма Флёри и контроль усвоения материала.

Цель работы состояла в разработке интерактивного тренажёра, который можно использовать в системе дистанционного обучения ИТМО для изучения и проверки навыков построения эйлерова цикла в неориентированном графе.

Для достижения поставленной цели были решены следующие задачи:
\begin{enumerate}
  \item проанализирована предметная область и сформулированы функциональные и нефункциональные требования к тренажёру;
  \item спроектированы структура приложения, модель данных и пользовательские сценарии;
  \item реализованы серверная и клиентская части приложения;
  \item разработан режим контролируемого тестирования с оцениванием качества решения;
  \item подготовлены материалы по развёртыванию и комплект отчётной документации в соответствии с требованиями методического пособия~\cite{markina2020,gost732}.
\end{enumerate}

В основной части отчёта подробно рассмотрено выполнение этапов 2--5 рабочего плана-графика: согласование функциональных требований, проектирование интерфейса и логики тренажёра, разработка и интеграция в LMS, тестирование и устранение ошибок. Для каждого этапа описаны не только полученные результаты, но и последовательность выполненных действий, принятые решения и выполненные доработки. В приложениях приведены материалы, подготовленные по итогам этих этапов.

\section{Согласование функциональных требований к тренажёру}

\subsection{Содержание этапов практики}

Работа в ходе практики была организована по утверждённому рабочему плану-графику. Для последующего анализа важно было зафиксировать не только конечный программный результат, но и логику прохождения этапов, поскольку именно она отражает процесс формирования требований, проектирования и программной реализации. Структура этапов приведена в таблице~\ref{tab:stages}.

\begin{table}[H]
\caption{Этапы выполнения практики по рабочему плану-графику}
\label{tab:stages}
\centering
\begin{tabular}{|c|p{4.7cm}|c|p{5.9cm}|}
\hline
\textbf{№} & \textbf{Наименование этапа} & \textbf{Период} & \textbf{Содержание работ} \\ \hline
1 & Инструктаж обучающегося & 1 & Ознакомление с требованиями охраны труда, техники безопасности, пожарной безопасности и правилами внутреннего распорядка. \\ \hline
2 & Обсуждение с куратором функциональных требований к интерактивному тренажёру по алгоритму Флёри & 3 & Определение формата визуализации графа, сценариев пошагового прохождения алгоритма, типов заданий и критериев оценки ответов. \\ \hline
3 & Проектирование интерфейса и логики тренажёра & 6 & Разработка структуры экрана, элементов управления, схемы взаимодействия модулей и прототипа пользовательского сценария. \\ \hline
4 & Разработка тренажёра и интеграция с LMS & 13 & Реализация серверной и клиентской частей, настройка режимов работы, параметризация запуска и подготовка конфигурации для среды дистанционного обучения. \\ \hline
5 & Тестирование и устранение ошибок & 13 & Подготовка набора проверок, выполнение функционального тестирования, анализ граничных случаев и исправление обнаруженных недостатков. \\ \hline
6 & Оформление отчётных документов в соответствии с требованиями & 15 & Подготовка отчёта, приложений и сопроводительных материалов по итогам практики. \\ \hline
7 & Получение отзыва руководителя практики & 7 & Передача результатов практики руководителю и получение итогового отзыва с оценкой. \\ \hline
\end{tabular}
\end{table}

В основной части далее подробно рассматриваются этапы 2--5, поскольку именно на них выполнялась содержательная работа по созданию тренажёра. Этапы 1, 6 и 7 носят сопроводительный характер: они обеспечивают формальное начало практики, оформление результатов и итоговую оценку.

\subsection{Выполнение этапа согласования требований}

Работа над проектом началась не с написания кода, а с обсуждения задачи с куратором практики. На этом этапе требовалось перевести общую тему практики в набор конкретных требований к программному средству. Обсуждение было сосредоточено на том, каким должен быть тренажёр с точки зрения учебного сценария, какие действия должен выполнять пользователь и какие результаты должны фиксироваться системой.

В ходе согласования были последовательно рассмотрены следующие вопросы:
\begin{itemize}
  \item как должен быть представлен граф на экране и какие действия пользователь должен выполнять вручную;
  \item должен ли тренажёр только демонстрировать готовое решение или также проверять действия обучающегося;
  \item какие режимы работы необходимы для учебного и контрольного использования;
  \item какие ограничения и критерии оценивания следует заложить в тестовый режим;
  \item каким образом подготовленный тренажёр будет запускаться в среде дистанционного обучения.
\end{itemize}

Первоначально задача могла быть сведена к обычной визуализации алгоритма Флёри, однако после обсуждения стало ясно, что этого недостаточно для целей практики. Было решено, что система должна поддерживать не один, а несколько сценариев использования: самостоятельное построение и редактирование графа, пошаговую демонстрацию алгоритма, автоматическое выполнение и отдельный режим контролируемого тестирования. Такое уточнение повлияло на все последующие этапы: уже на стадии проектирования потребовалось учитывать как обучающую, так и оценочную функции тренажёра.

Отдельным результатом этапа стало уточнение состава проверяемых действий пользователя. В тестовом режиме необходимо было оценивать не только факт построения эйлерова цикла, но и корректность выбора рёбер в соответствии с правилом Флёри. Это решение усложнило постановку задачи, поскольку потребовало реализовать анализ пользовательского ответа на уровне отдельных шагов, а не только на уровне конечного маршрута.

\subsection{Теоретическая основа и конкретизация постановки задачи}

После согласования общих требований была уточнена математическая основа тренажёра. Эйлеровым циклом называется замкнутый маршрут в графе, проходящий по каждому ребру ровно один раз. Для существования эйлерова цикла в неориентированном графе необходимо, чтобы граф был связным по вершинам ненулевой степени и степени всех его вершин были чётными~\cite{bondy2008,harary1969}. В отчёте это условие записывается следующим образом:

\begin{equation}
\forall v \in V \quad \deg(v) \equiv 0 \pmod{2},
\label{eq:euler}
\end{equation}
где $V$ -- множество вершин графа, $\deg(v)$ -- степень вершины $v$.

Далее была проанализирована специфика самого алгоритма Флёри. Он строит эйлеров цикл поэтапно: на каждом шаге из текущей вершины выбирается ребро, которое не является мостом в остаточном графе, если существует альтернативный выбор~\cite{wilson2010,west2001}. Из этого следовало, что для корректного тренажёра требуется не только вычислять очередной шаг, но и объяснять пользователю, почему данный шаг допустим или вынужден.

По итогам этапа 2 задача была сформулирована в рабочем виде как разработка веб-приложения для системы дистанционного обучения ИТМО, которое:
\begin{itemize}
  \item визуализирует неориентированный граф и ход алгоритма;
  \item позволяет пользователю самостоятельно задавать или генерировать эйлеров граф;
  \item поддерживает пошаговое и автоматическое выполнение алгоритма Флёри;
  \item содержит отдельный режим тестирования с автоматической оценкой ответа;
  \item может быть воспроизводимо запущено в стандартной среде LMS.
\end{itemize}

Именно в таком виде постановка задачи была передана на следующий этап, где требовалось уже не обсуждать требования, а проектировать конкретную структуру интерфейса и внутренней логики системы.

\section{Проектирование интерактивного тренажёра}

\subsection{Проработка пользовательских сценариев и требований}

После согласования общей постановки задачи работа была перенесена на уровень проектирования. На этом этапе требовалось превратить словесно сформулированные ожидания куратора в структуру экранов, набор элементов управления и последовательность пользовательских действий. Проектирование начиналось с выделения основных сценариев использования: редактирование графа, пошаговое прохождение алгоритма, автоматическая демонстрация и контролируемое тестирование.

Сначала был зафиксирован минимальный набор действий, который пользователь должен выполнять без переключения между разными страницами. Это решение было принципиальным: поскольку тренажёр предназначен для LMS, интерфейс должен оставаться компактным, а все основные операции должны быть доступны из одного экрана. В результате была выбрана двухпанельная композиция: слева размещается рабочая область с графом и результатами, справа -- элементы управления, настройки и блок тестирования.

Следующим шагом стало уточнение функциональных требований к каждому сценарию. Сводка ключевых требований, сформированная на этапе проектирования, приведена в таблице~\ref{tab:requirements}.

\begin{table}[H]
\caption{Ключевые требования к тренажёру}
\label{tab:requirements}
\centering
\begin{tabular}{|p{3.5cm}|p{8.7cm}|}
\hline
\textbf{Группа требований} & \textbf{Содержание} \\ \hline
Работа с графом & Добавление и удаление вершин и рёбер, очистка графа, генерация случайного эйлерова графа, выбор стартовой вершины. \\ \hline
Учебный режим & Пошаговый запуск алгоритма, автоматическое выполнение, журнал действий, визуальная подсветка стартовой и текущей вершины, отображение пройденных рёбер и мостов. \\ \hline
Тестовый режим & Генерация задания, ввод цикла Эйлера для каждой вершины, ограничение времени и числа попыток, автоматическая проверка и выставление баллов. \\ \hline
Надёжность вычислений & Валидация входных данных, проверка условий эйлеровости, защита от некорректных имён вершин и несуществующих рёбер. \\ \hline
Интеграция & Запуск в контейнерной инфраструктуре, настройка через переменные окружения, подготовка конфигурации для локальной разработки и production-сборки. \\ \hline
Педагогическая пригодность & Простая визуальная схема, минимальное число действий до получения результата, наглядная обратная связь при ошибках пользователя. \\ \hline
\end{tabular}
\end{table}

Таблица требований использовалась не как итоговая справка, а как рабочий инструмент проектирования. По ней последовательно проверялось, не потеряна ли какая-либо функция при переходе от одного проектного решения к другому. Например, требование по отображению мостов повлияло на состав визуальных состояний рёбер, а требование по контролируемому тестированию -- на появление отдельного механизма анализа ответов и системы начисления баллов.

\subsection{Проектирование архитектуры и логики взаимодействия}

После фиксации пользовательских сценариев был выполнен переход к проектированию внутренней структуры системы. На этом шаге нужно было определить, какие части логики следует реализовать на сервере, какие -- в браузере, и как организовать обмен данными между ними. Ключевым решением стало разделение системы на клиентский уровень, серверное Flask-приложение и отдельный алгоритмический слой. Такое разделение позволило не смешивать визуализацию с вычислительной логикой и обеспечило повторное использование алгоритмических функций в разных режимах работы.

Структура разработанного тренажёра представлена на рисунке~\ref{fig:architecture}.

\begin{figure}[H]
\centering
\fbox{%
\begin{minipage}{0.88\textwidth}
\centering
\textbf{Пользовательский интерфейс}\\
HTML, CSS, JavaScript, Cytoscape.js\\[4pt]
$\Downarrow$ HTTP/JSON\\[4pt]
\textbf{Flask-приложение}\\
маршруты API, валидация данных, переключение режимов training/testing\\[4pt]
$\Downarrow$ вызовы модулей\\[4pt]
\textbf{Алгоритмический слой}\\
модель графа, проверка эйлеровости, реализация алгоритма Флёри,\\
генерация графов и анализ пользовательских ответов\\[4pt]
$\Downarrow$ конфигурация запуска\\[4pt]
\textbf{Контейнерная среда}\\
Dockerfile, docker-compose, переменные окружения
\end{minipage}%
}
\caption{Структура интерактивного тренажёра}
\label{fig:architecture}
\end{figure}

В процессе проектирования были выделены следующие основные сущности серверной части:
\begin{itemize}
  \item модель графа, включающая список вершин и рёбер, операции модификации и базовые методы анализа;
  \item модуль реализации алгоритма Флёри, отвечающий за инициализацию, пошаговое выполнение и формирование текущего состояния;
  \item функции анализа пользовательского решения, проверяющие корректность пути и соответствие правилу Флёри;
  \item слой HTTP-маршрутов, предоставляющий интерфейс для клиентского приложения.
\end{itemize}

Выбор такой декомпозиции был связан с практическими соображениями. Если бы логика проверки и алгоритм Флёри находились только в клиентской части, было бы труднее обеспечивать единообразие поведения в учебном и тестовом режимах. Перенос вычислительной части на сервер позволил централизовать валидацию, упростить проверку ответов и уменьшить зависимость от состояния пользовательского интерфейса. Подобная архитектура соответствует типовой практике разработки веб-приложений на Python и Flask~\cite{python3docs,flaskdocs}.

\subsection{Уточнение модели данных и сценариев работы}

Отдельной задачей этапа проектирования стало определение формата представления графа. Необходимо было выбрать такую модель данных, которая, с одной стороны, позволяла бы выполнять алгоритмическую обработку на сервере, а с другой -- без лишних преобразований передавалась бы в браузер для визуализации. По этой причине была выбрана простая структура: список идентификаторов вершин и список рёбер вида <<источник -- цель>>. Такой формат удобно сериализуется в JSON и напрямую передаётся в клиентскую часть, где используется библиотека Cytoscape.js~\cite{cytoscapejs}.

После выбора модели данных был уточнён состав сценариев работы системы.

Первый сценарий -- учебный. Пользователь самостоятельно задаёт или генерирует граф, выбирает стартовую вершину и запускает алгоритм по шагам или автоматически. На этапе проектирования для этого сценария специально были заложены журнал действий, таблица маршрута и визуальная подсветка состояния графа. Эти элементы появились не как декоративные части интерфейса, а как ответ на требование сделать процесс прохождения алгоритма объяснимым.

Второй сценарий -- контрольный. Здесь проектирование было сложнее, поскольку требовалось совместить автоматическую проверку с понятной для пользователя обратной связью. Было решено, что система будет генерировать граф, принимать ответы по каждой стартовой вершине и после проверки возвращать не только итоговый балл, но и частные показатели качества решения: число корректно использованных рёбер, число замкнутых циклов и число нарушений правила Флёри. Наличие двух сценариев было зафиксировано как обязательное проектное решение, поскольку именно оно делало тренажёр одновременно учебным и контрольным инструментом.

\section{Разработка тренажёра и интеграция с LMS}

\subsection{Последовательность реализации серверной части}

Этап разработки был начат с серверной части, поскольку именно здесь сосредоточена алгоритмическая логика, определяющая корректность всего тренажёра. Реализация велась поэтапно: сначала были созданы базовые операции над графом, затем проверка условий эйлеровости, после этого -- механизм пошагового выполнения алгоритма, и только затем поверх этих сущностей был организован слой HTTP-маршрутов.

Серверная часть реализована на языке Python с использованием микрофреймворка Flask~\cite{python3docs,flaskdocs}. В модуле \texttt{app.py} инициализируется приложение, настраиваются режимы работы и описываются HTTP-маршруты. Выбор Flask был обусловлен необходимостью быстро построить компактное веб-приложение с явным разделением маршрутов, входных данных и режимов работы. Переменная окружения \texttt{APP\_MODE} определяет, запускается ли приложение в режиме обучения или тестирования, а дополнительные параметры задают число попыток, максимальный балл, временной лимит и диапазон размера графа.

Алгоритмическая логика вынесена в модуль \texttt{algorithm.py}. Ключевым элементом является класс \texttt{Graph}, в котором реализованы методы добавления и удаления вершин и рёбер, вычисления степени вершины, проверки достижимости и клонирования графа. Отдельная функция \texttt{check\_eulerian\_graph} проверяет выполнение условий эйлеровости, сформулированных в~\eqref{eq:euler}, а также связность графа по вершинам ненулевой степени. Такой порядок разработки позволил сначала получить устойчивое ядро для операций над графом, а затем строить поверх него прикладную логику тренажёра.

После этого был реализован класс \texttt{FleuryAlgorithm}, инкапсулирующий состояние пошагового выполнения. Он хранит исходный граф, его рабочую копию, стартовую вершину, построенный путь и список уже использованных рёбер. Основная трудность на данном шаге заключалась в том, чтобы корректно различать вынужденные переходы и случаи, когда пользователь должен избегать выбора моста. Для решения этой задачи была реализована отдельная проверка мостов в остаточном графе, а состояние шага было расширено служебной информацией, необходимой для последующей визуализации и объяснения выбора ребра пользователю.

Заметным элементом этапа реализации стала функция генерации случайного эйлерова графа. Изначально рассматривался вариант случайного добавления рёбер, однако он приводил к большому числу некорректных или слишком простых графов. Поэтому был выбран более устойчивый подход: сначала строится простой цикл по всем вершинам, а затем случайным образом <<переключаются>> дополнительные циклы. Такая операция сохраняет чётность степеней вершин и при дополнительной проверке позволяет получать более разнообразные, но корректные эйлеровы графы.

\subsection{Разработка клиентской части}

После стабилизации серверной логики началась реализация клиентской части. Работа над интерфейсом велась итеративно: сначала был собран минимальный рабочий экран, после чего по результатам внутренней проверки уточнялись расположение элементов, состав кнопок управления и способ отображения промежуточных состояний алгоритма.

Клиентская часть состоит из HTML-шаблона, CSS-стилей и нескольких JavaScript-модулей. Разметка страницы задаёт двухпанельную структуру: слева располагается область графа, журнал и таблица текущего маршрута, справа -- редактор графа, элементы управления алгоритмом, блок тестирования, легенда и описание. Такая структура появилась как прямое продолжение этапа проектирования: она позволила сохранить на одном экране и область визуализации, и управляющие элементы, и результаты работы алгоритма.

Визуализация графа выполняется библиотекой Cytoscape.js~\cite{cytoscapejs}. Для вершин и рёбер были последовательно настроены стили, отражающие их текущее состояние: стартовая вершина, текущая вершина, вершины маршрута, уже пройденные рёбра, мосты и повторно использованные рёбра. Это решение было принципиальным для учебного режима, поскольку позволяло превратить сухой вычислительный результат в наглядную пошаговую демонстрацию.

Для снижения связности клиентский код был разделён на модули. Модуль \texttt{training.js} отвечает за работу элементов редактора графа и учебного режима: через него обрабатываются добавление и удаление вершин и рёбер, очистка графа, генерация случайного примера и запуск алгоритма. Модуль \texttt{algorithm.js} реализует взаимодействие с сервером для инициализации, выполнения шага и автоматического прохождения алгоритма, а также обновляет журнал и таблицу маршрута. Модуль \texttt{testing.js} отвечает за сценарий контролируемого тестирования и формирует поля ввода для ответов по каждой стартовой вершине. Такое разделение облегчило последовательную доработку отдельных сценариев без переписывания всего интерфейса.

\subsection{Разработка механизма оценивания и подготовка к интеграции}

В рамках этапа разработки отдельно была реализована система оценивания тестового режима. Итоговая относительная оценка попытки вычисляется по формуле

\begin{equation}
R = 0{,}60 \cdot r_{\text{cov}} + 0{,}15 \cdot r_{\text{cycle}} + 0{,}25 \cdot r_{\text{fleury}},
\label{eq:score}
\end{equation}
где $r_{\text{cov}}$ -- доля корректно использованных рёбер, $r_{\text{cycle}}$ -- доля правильно замкнутых циклов, $r_{\text{fleury}}$ -- доля шагов, удовлетворяющих правилу Флёри.

Весовые коэффициенты в формуле~\eqref{eq:score} были подобраны в соответствии с согласованными критериями оценивания. Наибольший вклад в итоговую оценку даёт корректное прохождение по всем рёбрам, поскольку именно это отражает достижение основной учебной цели. Отдельно учитываются факт замыкания цикла и соблюдение правила выбора немостового ребра. После вычисления относительной оценки она преобразуется в дискретный балл, используемый системой дистанционного обучения.

Параллельно с реализацией логики оценивания велась подготовка приложения к интеграции в LMS. Для обеспечения воспроизводимости запуска были подготовлены файлы \texttt{Dockerfile}, \texttt{dev.yml}, \texttt{docker-compose.yml} и \texttt{spec.yml}. Конфигурация локальной разработки публикует сервис на порт \texttt{3000} и включает режим отладки, тогда как production-конфигурация задаёт параметры контейнера, лимиты ресурсов, сетевое окружение и набор переменных окружения для LMS~\cite{dockerdocs}. Это позволило завершить этап разработки не только рабочим кодом, но и подготовленной схемой запуска в учебной инфраструктуре.

\section{Тестирование и устранение ошибок}

\subsection{Подготовка сценариев проверки}

После завершения основной реализации был сформирован набор функциональных проверок. Работа на этом этапе велась от пользовательских сценариев к конкретным тестовым случаям: сначала были выделены ключевые действия пользователя, затем для каждого из них были определены ожидаемые результаты и возможные граничные состояния. Сводка подготовленных сценариев приведена в таблице~\ref{tab:tests}.

\begin{table}[H]
\caption{Основные сценарии функциональной проверки}
\label{tab:tests}
\centering
\begin{tabular}{|p{4.5cm}|p{7.7cm}|}
\hline
\textbf{Проверяемый сценарий} & \textbf{Ожидаемый результат} \\ \hline
Создание и редактирование графа & Вершины и рёбра корректно добавляются, удаляются и отображаются в интерфейсе. \\ \hline
Проверка неэйлерова графа & Алгоритм не запускается, пользователю возвращается объясняющее сообщение о причине отказа. \\ \hline
Пошаговое выполнение алгоритма & На каждом шаге корректно обновляются журнал, таблица маршрута и визуальное состояние графа. \\ \hline
Автоматический запуск & Алгоритм завершает построение эйлерова цикла без потери согласованности состояния интерфейса. \\ \hline
Режим тестирования & Ответы пользователя валидируются, ограничения по времени и попыткам соблюдаются, итоговый балл вычисляется корректно. \\ \hline
Повторный запуск в контейнере & Приложение поднимается через конфигурацию локальной разработки и production-описание без ручной правки кода. \\ \hline
\end{tabular}
\end{table}

Подготовленный набор проверок использовался не только для подтверждения корректности работы, но и как средство поиска мест, требующих доработки. Особое внимание было уделено граничным случаям: пустому графу, графу с нечётными степенями вершин, повторному прохождению по одному и тому же ребру, вводу несуществующих вершин, нарушению структуры ответа и завершению теста по истечении времени.

\subsection{Выявленные проблемы и выполненные исправления}

Тестирование проводилось итеративно. После каждого найденного недостатка выполнялось исправление, затем соответствующий сценарий запускался повторно, а смежные функции дополнительно проверялись на отсутствие регрессий. Такой подход оказался особенно важным для случаев, в которых клиентская и серверная части влияют друг на друга.

По результатам тестирования были выполнены следующие группы исправлений:
\begin{itemize}
  \item усилена валидация входных данных, чтобы алгоритм не запускался для пустых, неэйлеровых и некорректно заданных графов;
  \item уточнено обновление клиентского состояния после каждого шага алгоритма, чтобы журнал, таблица маршрута и визуальное состояние графа оставались согласованными;
  \item доработана логика выбора рёбер в пошаговом режиме, чтобы вынужденные переходы и мосты корректно определялись и отражались в интерфейсе;
  \item проверены ограничения тестового режима по времени, числу попыток и формату ответов пользователя;
  \item подтверждена воспроизводимость контейнерного запуска без ручного изменения исходного кода и конфигурации.
\end{itemize}

После внесения исправлений подготовленные сценарии были выполнены повторно. Повторная проверка показала, что ключевые пользовательские сценарии работают стабильно, а граничные случаи обрабатываются явно, без скрытых ошибок и неоднозначных состояний интерфейса. Тем самым этап 5 завершился не только фиксацией качества результата, но и реальным циклом поиска, анализа и устранения обнаруженных проблем.

\structuralsection{Заключение}

В ходе производственной, технологической (проектно-технологической) практики была выполнена разработка интерактивного тренажёра <<Алгоритм Флёри: построение эйлерова цикла в неориентированном графе>> для системы дистанционного обучения ИТМО.

В ходе этапа согласования требований были уточнены учебные и технические ожидания к тренажёру, определены ключевые сценарии использования и зафиксированы критерии оценивания ответов пользователя. На этапе проектирования были разработаны структура интерфейса, архитектура взаимодействия клиентской, серверной и алгоритмической частей, а также модель представления графа и пользовательских действий. На этапе разработки были последовательно реализованы серверная логика, клиентский интерфейс, система оценивания и контейнерная конфигурация для интеграции в LMS. На этапе тестирования был сформирован набор сценариев проверки, выявлены проблемные случаи и выполнены исправления, обеспечившие устойчивую работу системы.

Практический результат работы представляет собой самодостаточное учебное веб-приложение, пригодное для демонстрации алгоритма Флёри, самостоятельной тренировки и текущего контроля знаний. После завершения технических этапов были также подготовлены отчётные материалы в соответствии с требованиями методического пособия. Поставленная цель достигнута, задачи практики выполнены в полном объёме.

\structuralsection{Список использованных источников}

\vspace{0.35\baselineskip}
\enlargethispage{3\baselineskip}
\renewcommand{\refname}{}
\begin{thebibliography}{99}
\setlength{\topsep}{0pt}
\setlength{\itemsep}{0pt}
\setlength{\parskip}{0pt}
\bibitem{markina2020}
Маркина Т.~А., Пенской А.~В., Штенников Д.~Г., Авксентьева Е.~Ю., Ильина А.~Г. Производственная практика магистрантов: организация и проведение. СПб.: Университет ИТМО, 2020. 50~с.

\bibitem{gost732}
ГОСТ 7.32-2017. Система стандартов по информации, библиотечному и издательскому делу. Отчёт о научно-исследовательской работе. Структура и правила оформления. М.: Стандартинформ, 2017. 32~с.

\bibitem{wilson2010}
Wilson R.~J. Introduction to Graph Theory. 5th ed. Harlow: Pearson, 2010. 648~p.

\bibitem{west2001}
West D.~B. Introduction to Graph Theory. 2nd ed. Upper Saddle River: Prentice Hall, 2001. 588~p.

\bibitem{bondy2008}
Bondy J.~A., Murty U.~S.~R. Graph Theory. New York: Springer, 2008. 651~p.

\bibitem{harary1969}
Diestel R. Graph Theory. 5th ed. Berlin: Springer, 2017. 428~p.

\bibitem{python3docs}
Python Software Foundation. Python 3 Documentation [Электронный ресурс]. Режим доступа: \url{https://docs.python.org/3/}, свободный (дата обращения: 04.04.2026).

\bibitem{flaskdocs}
Pallets. Flask Documentation [Электронный ресурс]. Режим доступа: \url{https://flask.palletsprojects.com/}, свободный (дата обращения: 04.04.2026).

\bibitem{cytoscapejs}
Cytoscape.js Documentation [Электронный ресурс]. Режим доступа: \url{https://js.cytoscape.org/}, свободный (дата обращения: 04.04.2026).

\bibitem{mdnhtml}
WHATWG. HTML Living Standard [Электронный ресурс]. Режим доступа: \url{https://html.spec.whatwg.org/}, свободный (дата обращения: 04.04.2026).

\bibitem{mdncss}
W3C. CSS Snapshot [Электронный ресурс]. Режим доступа: \url{https://www.w3.org/TR/css/}, свободный (дата обращения: 04.04.2026).

\bibitem{mdnjs}
ECMA International. ECMAScript Language Specification [Электронный ресурс]. Режим доступа: \url{https://tc39.es/ecma262/}, свободный (дата обращения: 04.04.2026).

\bibitem{dockerdocs}
Docker Docs. Compose file reference [Электронный ресурс]. Режим доступа: \url{https://docs.docker.com/reference/compose-file/}, свободный (дата обращения: 04.04.2026).
\end{thebibliography}

\appendixsection{А}{Материалы этапов согласования требований и проектирования}

В таблице~\ref{tab:use-cases} приведены основные пользовательские сценарии, сформированные по результатам этапа анализа требований.

\begin{table}[H]
\caption{Основные пользовательские сценарии тренажёра}
\label{tab:use-cases}
\centering
\begin{tabular}{|p{3.8cm}|p{8.4cm}|}
\hline
\textbf{Сценарий} & \textbf{Краткое описание} \\ \hline
Редактирование графа & Пользователь создаёт вершины и рёбра, удаляет ошибочные элементы, очищает рабочее поле или генерирует случайный пример. \\ \hline
Пошаговое обучение & Пользователь выбирает стартовую вершину и последовательно запускает шаги алгоритма, наблюдая журнал и маршрут. \\ \hline
Автоматическая демонстрация & Система полностью выполняет алгоритм и показывает готовый эйлеров цикл. \\ \hline
Контролируемое тестирование & Пользователь вводит эйлеровы циклы, а система проверяет корректность рёбер, замыкание циклов и соблюдение правила Флёри. \\ \hline
Публикация результата в LMS & После завершения теста система возвращает итоговый балл и формирует данные для внешней фиксации результата. \\ \hline
\end{tabular}
\end{table}

На рисунке~\ref{fig:appendix-arch} приведена укрупнённая схема взаимодействия компонентов, использованная на этапе проектирования.

\begin{figure}[H]
\centering
\fbox{%
\begin{minipage}{0.88\textwidth}
\centering
Пользователь $\rightarrow$ интерфейс браузера $\rightarrow$ HTTP API $\rightarrow$\\
модуль алгоритма и проверки $\rightarrow$ результат и визуальная обратная связь
\end{minipage}%
}
\caption{Укрупнённая схема взаимодействия компонентов}
\label{fig:appendix-arch}
\end{figure}

\appendixsection{Б}{Материалы этапа разработки}

Ниже приведён сокращённый фрагмент функции генерации случайного эйлерова графа, отражающий принцип сохранения чётности степеней вершин.

\begin{lstlisting}[language=Python, caption={Фрагмент функции генерации эйлерова графа}]
def generate_eulerian_graph(min_nodes=6, max_nodes=9):
    node_count = random.randint(min_nodes, max_nodes)
    nodes = [chr(ord('A') + i) for i in range(node_count)]

    edge_keys = set()
    for i in range(node_count):
        source = nodes[i]
        target = nodes[(i + 1) % node_count]
        edge_keys.add(Graph.edge_key(source, target))

    target_ops = random.randint(max(2, node_count // 2), max(4, node_count * 2))
    attempts_left = target_ops * 8

    while target_ops > 0 and attempts_left > 0:
        attempts_left -= 1
        cycle_len = random.randint(3, node_count)
        cycle_nodes = random.sample(nodes, cycle_len)
        cycle_edges = [
            Graph.edge_key(cycle_nodes[i], cycle_nodes[(i + 1) % cycle_len])
            for i in range(cycle_len)
        ]
\end{lstlisting}

Следующий фрагмент отражает центральную часть реализации алгоритма Флёри: на каждом шаге из текущей вершины выбирается допустимое ребро, причём при наличии альтернативы предпочтение отдаётся ребру, не являющемуся мостом в остаточном графе.

\begin{lstlisting}[language=Python, caption={Фрагмент пошагового выполнения алгоритма Флёри}]
def step(self) -> Dict:
    if self.finished:
        return self.get_state()

    neighbors = self.graph.get_neighbors(self.current_node)

    if not neighbors:
        self.finished = True
        if self.graph.edge_count() == 0 and len(self.path) > 1 and self.path[0] == self.path[-1]:
            self.valid = True
            self.message = 'Эйлеров цикл построен.'
        else:
            self.valid = False
            self.message = 'Цикл не построен: закончились доступные ребра.'
        return self.get_state()

    non_bridge_neighbors = [
        neighbor for neighbor in neighbors
        if not self._is_bridge(self.current_node, neighbor)
    ]

    forced = len(neighbors) == 1 or len(non_bridge_neighbors) == 0
    next_node = sorted(neighbors)[0] if forced else sorted(non_bridge_neighbors)[0]
    is_bridge = self._is_bridge(self.current_node, next_node)

    self.graph.remove_edge(self.current_node, next_node)
\end{lstlisting}

Следующий фрагмент иллюстрирует вычисление качества пользовательского решения в тестовом режиме.

\begin{lstlisting}[language=Python, caption={Фрагмент вычисления оценки решения}]
SCORE_WEIGHTS = {
    'edge_coverage': 0.60,
    'cycle_complete': 0.15,
    'fleury_rule': 0.25,
}

def _calculate_attempt_score(cycle_summary):
    coverage_ratio = cycle_summary.get('coverage_ratio', 0.0)
    complete_ratio = cycle_summary.get('complete_ratio', 0.0)
    fleury_ratio = cycle_summary.get('fleury_ratio', 0.0)

    weighted_ratio = (
        coverage_ratio * SCORE_WEIGHTS['edge_coverage']
        + complete_ratio * SCORE_WEIGHTS['cycle_complete']
        + fleury_ratio * SCORE_WEIGHTS['fleury_rule']
    )

    attempt_score = int(TEST_TOTAL_SCORE * weighted_ratio)
    return max(0, attempt_score)
\end{lstlisting}

\appendixsection{В}{Материалы этапа интеграции с LMS}

Для локальной разработки и интеграции в LMS были подготовлены разные конфигурации запуска. В листинге~\ref{lst:devcfg} показан сокращённый фрагмент локальной конфигурации.

\begin{lstlisting}[language={}, caption={Фрагмент локальной конфигурации запуска}, label={lst:devcfg}]
services:
  task:
    build:
      context: ./src
    ports:
      - "3000:80"
    environment:
      - APP_MODE=${APP_MODE:-training}
      - TEST_MAX_ATTEMPTS=3
      - TEST_TOTAL_SCORE=5
      - TEST_MIN_SCORE=3
      - TEST_MAX_SCORE=5
      - TEST_TIME_LIMIT_MINUTES=60
      - DEBUG=1
\end{lstlisting}

В production-конфигурации, напротив, задаются лимиты ресурсов, параметры сетевого окружения и переменные, необходимые для взаимодействия с внешней инфраструктурой системы дистанционного обучения.

\appendixsection{Г}{Материалы этапа тестирования}

В таблице~\ref{tab:docs} приведены характерные ситуации, проверявшиеся на этапе тестирования, и результаты выполненных исправлений.

\begin{table}[H]
\caption{Характерные результаты этапа тестирования}
\label{tab:docs}
\centering
\begin{tabular}{|p{6cm}|p{6cm}|}
\hline
\textbf{Проверяемая ситуация} & \textbf{Результат и выполненная доработка} \\ \hline
Пустой или неэйлеров граф & Добавлена явная проверка корректности графа до запуска алгоритма, пользователю выводится объясняющее сообщение. \\ \hline
Пошаговое выполнение алгоритма & Уточнена синхронизация состояния между сервером и интерфейсом: журнал, маршрут и подсветка обновляются согласованно. \\ \hline
Определение мостов и вынужденных ходов & Скорректирована логика отображения шага, чтобы в интерфейсе различались обычный выбор ребра и вынужденный переход. \\ \hline
Контролируемое тестирование & Проверены ограничения по времени, числу попыток и формату ввода; некорректные ответы обрабатываются без аварийных состояний. \\ \hline
Повторный контейнерный запуск & Подтверждено, что приложение стартует через подготовленные конфигурации без ручной правки исходного кода. \\ \hline
\end{tabular}
\end{table}

Таким образом, в приложениях зафиксированы материалы этапов 2--5: результаты согласования требований и проектирования, ключевые фрагменты реализации, конфигурация интеграции с LMS и итоги тестирования с выполненными исправлениями.

\end{document}
